% THEORY.tex — the theory behind the order book design.
% Compile: pdflatex THEORY.tex  (or upload to Overleaf — standard packages only)
\documentclass[11pt]{article}

\usepackage[a4paper,margin=2.4cm]{geometry}
\usepackage{xcolor}
\usepackage{titlesec}
\usepackage{enumitem}
\usepackage{fancyhdr}
\usepackage{mdframed}
\usepackage{listings}

\definecolor{accent}{RGB}{20,60,120}
\definecolor{buildbg}{RGB}{238,243,250}
\definecolor{codegray}{RGB}{245,245,245}

\titleformat{\section}{\Large\bfseries\color{accent}}{\thesection}{0.8em}{}
\titleformat{\subsection}{\large\bfseries}{\thesubsection}{0.7em}{}
\titleformat{\subsubsection}{\normalsize\bfseries}{\thesubsubsection}{0.6em}{}

\newmdenv[backgroundcolor=buildbg,linewidth=0pt,innerleftmargin=10pt,
          innerrightmargin=10pt,innertopmargin=8pt,innerbottommargin=8pt,
          skipabove=10pt,skipbelow=10pt]{notebox}

% compact verbatim-style block for diagrams (listings handles nesting safely)
\lstnewenvironment{diagram}
  {\lstset{basicstyle=\ttfamily\small,
           columns=fullflexible,
           keepspaces=true,
           backgroundcolor=\color{codegray},
           frame=none,
           xleftmargin=8pt,
           aboveskip=8pt,
           belowskip=8pt}}
  {}

\pagestyle{fancy}
\fancyhf{}
\fancyhead[R]{\small\itshape LOB project --- theory}
\fancyfoot[C]{\small\thepage}

\setlist{itemsep=2pt,topsep=4pt}
\setlength{\parskip}{4pt}

\title{\bfseries The Order Book Project:\\[4pt]
       \large The Theory Behind the Design}
\author{}
\date{}

\begin{document}
\maketitle
\thispagestyle{fancy}

\noindent A limit order book is a piece of exchange infrastructure: the thing
between ``an order arrives'' and ``a trade happens.'' This document collects
the theory that explains why the fast implementation is built the way it is
--- because in this domain, the theory \emph{is} the implementation. Nothing
here depends on a specific language or codebase; every idea is about memory,
complexity, and what hardware can answer natively.

%==================================================================
\section{What a book actually is}

\subsection{The wall of notes}

A farmer's market has a wall where buyers pin ``I'll buy at \$1.00, need 10''
and sellers pin ``selling at \$1.05, have 20.'' The \emph{book} is that wall,
kept sorted. The \emph{matching engine} is the clerk who rips notes off when
they match. Two rules govern everything:

\begin{enumerate}
  \item \textbf{Price priority} --- a new buyer trades with the
        \emph{cheapest} seller first. Best price wins.
  \item \textbf{Time priority} --- at the same price, earliest note trades
        first. FIFO queue per level.
\end{enumerate}

``Price-time priority'' is the entire spec. The only verbs are
add, cancel, and modify; the only output is trade reports, always executed
at the \emph{resting} order's price. A real exchange is this loop, millions
of times per second.

\subsection{Correctness is defined by behavior, not internals}

Two implementations of the same observable contract are interchangeable by
definition. That gives two consequences worth internalizing:

\begin{itemize}
  \item \textbf{State machine view.} The book is a deterministic machine:
        events in, state transitions, trades out.
  \item \textbf{Differential testing.} The same tape through two
        implementations must yield identical trades. ``Is it correct?''
        becomes ``does it match the reference?'' --- falsifiable forever,
        for free.
\end{itemize}

Get it right first; every optimization afterward is continuously verified
against a known-good oracle.

%==================================================================
\section{The O(1) redesign}

\subsection{Memory is the cost, not compute}

The CPU is roughly 100x faster at arithmetic than at fetching from DRAM:

\begin{diagram}
executing one instruction on data in a register:   ~0.3 ns
L1 hit ~1 ns | L2 ~12 cy | L3 ~45 cy | DRAM ~200-300 cy (~100 ns)
\end{diagram}

The gap exists physically (DRAM is a separate chip, capacitors, distance)
and historically (CPU speed grew ~50\%/yr while DRAM latency improved
~7\%/yr --- the ``memory wall''; caches are the workaround). \textbf{One
cache miss $\approx$ 300 instructions} --- so op-count alone is the wrong
cost model.

\subsection{Dependent vs.\ independent misses}

Misses only serialize if they're \emph{dependent} --- when the next address
is hidden inside data you haven't fetched yet:

\begin{diagram}
independent (array):  known addresses -> fire all -> waits OVERLAP
dependent  (tree):    node -> next addr inside it -> serial, ~100ns each
\end{diagram}

The CPU keeps ten or more fetches in flight, but only for computable
addresses. Pointer chasing forbids overlap; the prefetcher can't help
either. \textbf{Tree and linked structures are the worst case:} maximum
dependent misses, zero parallelism.

\subsection{The address formula and its three conditions}

\begin{diagram}
f(i) = base + i * size     <- works only if:
  1. fixed-size elements   (constant stride; variable size => scan)
  2. stored INLINE         (an array of pointers is contiguous but still
                            chases; the object must BE at the address)
  3. back-to-back          (contiguous; else no formula exists)
\end{diagram}

Bonus: cache fetches in \textbf{64-byte lines} --- contiguity means each
miss carries many useful elements, not one. Overlap fixes \emph{latency};
lines fix \emph{fetch count}. You need both.

\subsection{Escaping the comparison model}

$\Omega(\log N)$ ordered access is a theorem \emph{only} for
comparison-based structures (``is $a < b$?'' yields one bit per step).
Prices are bounded integers --- an integer in a known range isn't just
orderable, \emph{it's an address}: subtract the base and index directly.

Same escape as radix sort beating $\Omega(N \log N)$. The honest limit: it
costs memory proportional to the universe, and it only works because prices
are bounded and dense --- unbounded keys put you back in the comparison
model. The tree's logarithmic access is \emph{unconditional}; that
generality is what the log factor buys.

\subsection{Bitmaps and bitscan: asking the hardware}

The level array is mostly empty --- we need ``which prices are live''
without scanning. One bit per price; finding the best ask is a single
count-trailing-zeros \emph{instruction}, not a loop: all 64 bits exist as
voltages simultaneously and a priority-encoder circuit answers in ~1 cycle.

\begin{diagram}
prices:  4900 4901 4903 4904
bitmap:    0    0    1    1       lowest set bit = best ask
                  ^
one word covers 64 prices; 10k prices = ~1.2KB = permanently L1-resident
bids scan the other direction (highest bit); "next nonempty" is the
same instruction applied word by word
\end{diagram}

Pattern worth naming: \emph{encode the problem as a question the hardware
answers natively} (chess bitboards, OS scheduler runqueues). The worst case
is universe-proportional word scans --- fixable with a hierarchical bitmap
(bits over words), the practical cousin of van Emde Boas layouts.

\subsection{Mutation is where trees bleed}

Lookups were the warmup; a book mostly \emph{writes}. A tree insert costs a
heap allocation plus a descent plus a rebalance --- writes to ancestors you
didn't ask about. And heap allocation is a global, lock-protected,
\textbf{unbounded} call: tens of nanoseconds usually, microseconds when the
heap is in a bad state. Amortized averages don't matter when the metric is
the tail. A flat design's whole mutation path is a few bounded writes to
computable addresses --- ``bounded tail'' means \emph{there is no code path
that can take milliseconds}.

\subsection{The pool, and the pardoned linked list}

\begin{itemize}
  \item \textbf{Pool/arena for orders}: one big array, allocation is a
        stack pop. No per-order heap traffic in the hot path; stable
        addresses mean identity-to-node pointers stay valid forever.
  \item \textbf{Intrusive FIFO per level}: orders link to each other;
        cancel is unlinking two pointers, constant time, no search. It is
        still a linked list --- dependent-chase memory --- but you only
        ever touch head and tail, so the chase is \emph{bounded at one
        hop}. The rule isn't ``never chase pointers,'' it's ``never chase
        unboundedly.''
  \item \textbf{Identity lookup} (``cancel order 47''): a flat hash map
        probes contiguous slots --- independent misses --- versus chained
        buckets that scatter. If ids are dense, direct indexing beats
        hashing entirely: the bounded-integer escape applied a second time.
\end{itemize}

\subsection{The composite is the real structure}

Two indexes over the same objects --- one by price, one by id --- must stay
in sync on every mutation. A fill that removes an order from its level but
leaves it in the id index creates a phantom; the reverse creates invisible
liquidity. Textbooks give you structures; \emph{keeping indexes consistent
under mutation is the actual engineering.}

%==================================================================
\section{Measuring like a professional}

\subsection{Percentiles are the truth; the mean is a summary}

Tail latency is the product. The 99.9th percentile exposes allocator slow
paths, rehashes, and rebalances that a mean would average away. A fair
benchmark controls warmup, allocator state, and seed --- and ``my fast
version was slower at X, here's why'' earns more credibility than any
headline number. You are measuring your code \emph{and} your harness; know
which is which.

\subsection{What the numbers can claim}

Matching $k$ levels costs $O(k)$ in fills --- inherent work, correctly
paid. What a fast book removes is per-operation overhead and unbounded
variance. The defensible claim is: \emph{``I exploited bounded dense
integer keys to leave the comparison model, and contiguous inline layout to
leave pointer chasing --- and here is what my design still cannot do.''}

%==================================================================
\section{Summary: the whole design in one breath}

Every choice in the fast book is the same two moves applied to different
data. \textbf{Move one}: turn keys into addresses --- prices index a flat
array, order ids index a hash map or dense table, occupancy lives in a
bitmap the CPU can scan in one instruction. \textbf{Move two}: make every
miss bounded and independent --- pool allocation instead of the heap,
intrusive lists bounded at one hop, sequential bitmap words the prefetcher
can chase for you.

The honest closing note is that none of this wins unconditionally. A small,
dense book keeps a balanced tree entirely cache-resident, and its best-node
lookup is genuinely constant time. The fast design's advantage is
\emph{scale and predictability}: as the book deepens and the price domain
widens, its costs stay flat while the tree's grow logarithmic and cold ---
and its worst case is a bounded scan rather than an allocator's bad day.
Knowing \emph{when} a structure wins, and being able to say why, is the
skill the structure was built to teach.

\end{document}
